\scanpage{119}
% archive sheet number: 119

\handwritten{2}. Supposons donné une situation $(Y,V,H,V')$ analogue à $(X,\handwritten{U},G,U')$, et un morphisme $f:X\handwritten{\longrightarrow}Y$ tel que $f(U)\handwritten{\longrightarrow}f(V)$, un morphisme $f':U'\handwritten{\longrightarrow}V'$ compatible avec $U\handwritten{\longrightarrow}V$ induit par $f$, un morphisme de groupes $\handwritten{\phi}:G\handwritten{\longrightarrow}H$ tel que \deleted{$\unsure{f'(gx')=}$} $f'$ soit compatible avec les actions de $G$ et $H$ via $\handwritten{\phi}$, \deleted{et}

Le cas le plus courant sera celui où on se donne un point géométrique $a$ de $U$, \inserted{d'où} un point géométrique $b$ \inserted{$=f(a)$} de $V$, où $U'$ et $V'$ sont respectivement des quotients des revêtements universels de $U$ et $V$ basés en $a$ et $b$, $G$ et $H$ étant respectivement des quotients de $\handwritten{\pi}_{1}(U,a)$ et $\handwritten{\pi}_{1}(V,b)$, l'homomorphisme $f':U'\handwritten{\longrightarrow}V'$ étant induit par les homomorphismes canoniques sur les revêtements universels en $a$ resp. $b$\deleted{\unsure{;|,} \unsure{un tel} $\unsure{f'}$} \unsure{\inserted{\deleted{ainsi que $\handwritten{\phi}$}}} \deleted{\unsure{est} \unsure{donc} déterminé de}, \deleted{façon unique $G$ et $H$ étant les g} $\handwritten{\phi}:G\handwritten{\longrightarrow}H$ étant induit par $f_{\handwritten{*}}:\handwritten{\pi}_{1}(U,a)\handwritten{\longrightarrow}\handwritten{\pi}_{1}(\unsure{\deleted{B}V|\deleted{V}B},b)$ ; de cette façon $f'$ et $\handwritten{\phi}$ sont déterminés sans ambiguïté par $f$, et leur existence signifie simplement que $f_{\handwritten{*}}:\handwritten{\pi}_{1}(U,a)\handwritten{\longrightarrow}\handwritten{\pi}_{1}(V,b)$ veut bien passer au quotient en un $\handwritten{\phi}:G\handwritten{\longrightarrow}H$ (i.e. applique un certain sous-groupe dans un certain sous-groupe). Considérons d'autre part le localisé strict $\handwritten{\overline}{Y}$ de $Y$ en le point géométrique $f(\handwritten{\bar}{x})$ de $Y$, d'où \deleted{\unsure{un}} $\handwritten{\bar}{V},\handwritten{\bar}{V}'$ comme dessus; supposons donné un prolongement $\handwritten{\overline}{Y'}$ de $\handwritten{\bar}{V}'$ en un schéma à groupe d'opérateurs $H$ sur $\handwritten{\bar}{Y}$, et supposons donné enfin un homomorphisme $\handwritten{\overline}{f'}:\handwritten{\overline}{X'}\handwritten{\longrightarrow}\handwritten{\overline}{Y'}$ qui prolonge l'homomorphisme \deleted{\unsure{ft|f'}} $\handwritten{\overline}{U'}\handwritten{\longrightarrow}\handwritten{\overline}{V'}$ induit par $f'$. Pour tout point $\handwritten{\xi}$ de $\handwritten{\overline}{X'}_{\handwritten{\bar}{x}}$, $\handwritten{\overline}{f'}(\handwritten{\xi})$ est alors un point\deleted{s} de $\handwritten{\overline}{Y'}_{\handwritten{\bar}{y}}$, et on peut considérer le groupe d'inertie correspondant $H_{\handwritten{\overline}{f'}(\handwritten{\xi})}$. On a alors la relation évidente
\[
\handwritten{\phi}(G_{\handwritten{\xi}})\handwritten{\subset}H_{\handwritten{\overline}{f'}(\handwritten{\xi})}\ .
\tag*{\handwritten{(2.1)}}
\]
Dans le cas où le\deleted{s} groupe $G$ resp. $H$ opère transitivement sur $\handwritten{\overline}{X'}_{\handwritten{\bar}{x}}$ resp. $\handwritten{\overline}{Y'}_{\handwritten{\overline}{y'}}$, on peut donc dire que $\handwritten{\phi}$ \inserted{applique tout groupe \deleted{de}} \deleted{transforme la classe de conjugaison de sous-groupes} d'inertie de $G$ rel. à $x$ \inserted{\handwritten{\uline{dans}} un groupe d\deleted{e}\handwritten{'}} \deleted{\unsure{en|à} la classe de conj. des sous-groupes} d'inertie de $H$ rel à $y=f(x)$.

\deleted{D'autre part} Il reste alors à expliciter l'homomorphisme $\handwritten{\overline}{f'}:\handwritten{\overline{X'}\longrightarrow\overline{Y'}}$.

\figure{\handwritten{\begin{tikzpicture}[x=0.006cm,y=-0.006cm]
\draw[gray,line width=0.3pt] (1664,451) .. controls (1622,413) and (1571,429) .. (1514,436) .. controls (1119,456) and (709,433) .. (276,428);
\draw[gray,line width=0.3pt] (258,438) .. controls (270,417) and (240,422) .. (228,451) .. controls (206,521) and (206,664) .. (202,870) .. controls (201,912) and (214,941) .. (218,979) .. controls (230,1065) and (230,1149) .. (211,1200) .. controls (236,1219) and (296,1206) .. (347,1208) .. controls (591,1211) and (732,1211) .. (901,1210) .. controls (950,1208) and (976,1188) .. (969,1152) .. controls (1209,1140) and (1415,1155) .. (1601,1137);
\draw[gray,line width=0.3pt,->] (211,1200) .. controls (186,1269) and (152,1426) .. (156,1546) .. controls (153,1742) and (178,1921) .. (203,1996) .. controls (212,2022) and (250,2009) .. (309,2019);
\draw[gray,line width=0.3pt] (63,532) -- (42,549) .. controls (33,559) and (36,583) .. (47,596);
\end{tikzpicture}}}
